% =====================================================================
%  MACROS
%  Símbolos que se repiten en todo el libro. Escribe SIEMPRE con macro,
%  nunca a mano: si un día cambias el convenio, cambias solo esta hoja.
% =====================================================================

% --- Generales --------------------------------------------------------
\usepackage{xspace}
\newcommand{\FUSAL}{\textsc{FUSAL}\xspace}
\newcommand{\etal}{\textit{et al.}\xspace}
\newcommand{\dif}{\mathrm{d}}
\newcommand{\deriv}[2]{\frac{\dif #1}{\dif #2}}
\newcommand{\pderiv}[2]{\frac{\partial #1}{\partial #2}}
\newcommand{\vect}[1]{\mathbf{#1}}
\newcommand{\prom}[1]{\overline{#1}}
% Texto guía dentro de los recuadros: qué hay que escribir ahí. Bórralo al redactar.
\newcommand{\redactar}[1]{\textcolor{fusalGris}{\itshape #1}}

% --- Dinámica vehicular ----------------------------------------------
\newcommand{\Fz}{F_z}                       % carga vertical
\newcommand{\Fy}{F_y}                       % fuerza lateral
\newcommand{\Fx}{F_x}                       % fuerza longitudinal
\newcommand{\Mz}{M_z}                       % par autoalineante
\newcommand{\slipa}{\alpha}                 % ángulo de deriva
\newcommand{\slipr}{\kappa}                 % deslizamiento longitudinal
\newcommand{\camber}{\gamma}                % caída
\newcommand{\muy}{\mu_y}
\newcommand{\mux}{\mu_x}
\newcommand{\ay}{a_y}
\newcommand{\ax}{a_x}
\newcommand{\CdG}{\mathrm{CdG}}             % centro de gravedad
\newcommand{\hcg}{h_{\CdG}}
\newcommand{\LLTD}{\mathrm{LLTD}}
\newcommand{\Kus}{K_{us}}                   % gradiente de subviraje
\newcommand{\batalla}{L}
\newcommand{\via}{t}
\newcommand{\Kroll}{K_{\phi}}               % rigidez a balanceo

% --- Neumático y su modelado (capítulos 8 y 9) ------------------------
\newcommand{\Calpha}{C_{\slipa}}            % rigidez de deriva
\newcommand{\Ckappa}{C_{\slipr}}            % rigidez de deslizamiento
\newcommand{\tneu}{t_p}                     % avance neumático
\newcommand{\pneu}{p_i}                     % presión de inflado
\newcommand{\Tneu}{T_n}                     % temperatura del neumático
\newcommand{\Fzo}{F_{z0}}                   % carga vertical de referencia
\newcommand{\dfz}{\mathrm{d}f_z}            % variación normalizada de carga
\newcommand{\lamu}{\lambda_{\mu}}           % factor de escala de agarre
% Coeficientes de la Magic Formula. OJO: \Cp ya es el coeficiente de
% presión aerodinámico, por eso estos llevan prefijo MF.
\newcommand{\MFB}{B}                        % factor de rigidez
\newcommand{\MFC}{C}                        % factor de forma
\newcommand{\MFD}{D}                        % valor de pico
\newcommand{\MFE}{E}                        % factor de curvatura
\newcommand{\MFSh}{S_h}                     % desplazamiento horizontal
\newcommand{\MFSv}{S_v}                     % desplazamiento vertical

% --- Suspensión -------------------------------------------------------
\newcommand{\MR}{\mathrm{MR}}               % relación de instalación
\newcommand{\Kw}{K_w}                       % rigidez en rueda
\newcommand{\Ks}{K_s}                       % rigidez del muelle
\newcommand{\fr}{f_r}                       % frecuencia de suspensión

% --- Aerodinámica -----------------------------------------------------
\newcommand{\Cl}{C_L}
\newcommand{\Cd}{C_D}
\newcommand{\Cp}{C_p}
\newcommand{\ClA}{C_L A}
\newcommand{\CdA}{C_D A}
\newcommand{\qdin}{q_{\infty}}              % presión dinámica
\newcommand{\Vinf}{V_{\infty}}
\newcommand{\Rey}{\mathit{Re}}
\newcommand{\CoP}{\mathrm{CoP}}             % centro de presiones

% --- Transferencia de calor ------------------------------------------
\newcommand{\UA}{UA}
\newcommand{\NTU}{\mathit{NTU}}
\newcommand{\efec}{\varepsilon}             % eficacia del intercambiador
\newcommand{\Nus}{\mathit{Nu}}
\renewcommand{\Pr}{\mathit{Pr}}   % redefine el \Pr de LaTeX (probabilidad), aquí es Prandtl
\newcommand{\Stan}{\mathit{St}}
\newcommand{\jcol}{j}                       % factor de Colburn
\newcommand{\ffan}{f}                       % factor de fricción
\newcommand{\mdot}{\dot{m}}
\newcommand{\Qdot}{\dot{Q}}
\newcommand{\Dh}{D_h}                       % diámetro hidráulico
\newcommand{\LMTD}{\Delta T_{\mathrm{ml}}}

% --- Transmisión ------------------------------------------------------
\newcommand{\itot}{i_{\mathrm{tot}}}        % relación total
\newcommand{\Tmot}{T_{\mathrm{mot}}}
\newcommand{\rdin}{r_{\mathrm{din}}}        % radio dinámico
